\section{Main results}\label{sec:main}

\subsection{Existence and characterization of the threshold}

Our first result establishes that the misalignment sensitivity threshold is
well-defined.

\begin{theorem}[Existence]\label{thm:existence}
Let $\mathcal{M}$ be a finite trading MDP with softmax policy parameterization
at temperature $T > 0$. For any perturbation direction $\Phi$ with
$\lVert \Phi \rVert_\infty = 1$ and significance level $\delta > 0$, the
misalignment sensitivity threshold $\tau(\delta)$ exists and satisfies
$0 < \tau(\delta) < \infty$.
\end{theorem}

\begin{proof}
We establish each bound separately.

\emph{Lower bound} ($\tau > 0$). The Q-function $Q_{R_\varepsilon}(s,a)$ is
the unique fixed point of the Bellman operator with reward
$R_\varepsilon = R_{\mathrm{proxy}} + \varepsilon \Phi$. Since $\Phi$ is
bounded, the Bellman operator depends continuously on $\varepsilon$, and by
the contraction mapping theorem the fixed point $Q_{R_\varepsilon}$ is
continuous in $\varepsilon$. The softmax map $Q \mapsto \pi^T$ is smooth for
$T > 0$ (it is a composition of exponentials and division by a positive
sum). Therefore $\pi_\varepsilon = \pi_{R_\varepsilon}^T$ is continuous in
$\varepsilon$, and hence $D(\varepsilon) = D_{\mathrm{KL}}(\pi_\varepsilon
\| \pi_0)$ is continuous in $\varepsilon$ with $D(0) = 0 < \delta$. By
continuity, there exists $\varepsilon_0 > 0$ such that $D(\varepsilon) <
\delta$ for all $\varepsilon \in [0, \varepsilon_0)$, whence
$\tau \geq \varepsilon_0 > 0$.

\emph{Upper bound} ($\tau < \infty$). For $\varepsilon$ sufficiently large,
the term $\varepsilon \Phi$ dominates $R_{\mathrm{proxy}}$ in the Bellman
equations, so the Q-function of $R_\varepsilon$ is asymptotically determined
by $\Phi$ alone. Concretely, for any state $s$ where
$\operatorname{argmax}_a \Phi(s,a)$ is a singleton not equal to
$\operatorname{argmax}_a Q_{R_{\mathrm{proxy}}}(s,a)$, the softmax policy
$\pi_\varepsilon(\cdot \mid s)$ concentrates on
$\operatorname{argmax}_a \Phi(s,a)$ as $\varepsilon \to \infty$, producing a
limiting policy $\pi_\infty$ with $D_{\mathrm{KL}}(\pi_\infty \| \pi_0) > 0$.
For generic $\Phi$ (which is dense in the unit ball of $\ell^\infty$), such
states exist. Therefore there exists $\varepsilon_1 < \infty$ with
$D(\varepsilon_1) > \delta$, giving $\tau \leq \varepsilon_1 < \infty$.
\end{proof}

The following lemma computes the first-order response of the softmax policy
to reward perturbations.

\begin{lemma}[Policy sensitivity formula]\label{lem:sensitivity}
For a softmax policy $\pi^T$ at temperature $T > 0$ with Q-function $Q$, the
policy gradient sensitivity is
\begin{equation}\label{eq:sensitivity}
G(s,a) = \frac{1}{T}\, \pi^T(a \mid s) \Bigl(\widetilde{\Phi}(s,a) -
  \mathbb{E}_{\pi^T(\cdot \mid s)}\bigl[\widetilde{\Phi}(s,\cdot)\bigr]
  \Bigr),
\end{equation}
where $\widetilde{\Phi}(s,a) = \partial Q_{R_\varepsilon}(s,a)/\partial
\varepsilon \big|_{\varepsilon=0}$ is the Q-function sensitivity to the
perturbation.
\end{lemma}

\begin{proof}
The softmax policy is
\[
\pi_\varepsilon^T(a \mid s) = \frac{\exp\bigl(Q_{R_\varepsilon}(s,a)/T
  \bigr)}{\sum_{b \in \mathcal{A}} \exp\bigl(Q_{R_\varepsilon}(s,b)/T\bigr)}.
\]
This is the standard softmax function applied to the vector
$Q_{R_\varepsilon}(s,\cdot)/T$. The derivative of the softmax output $z_i =
e^{x_i}/\sum_j e^{x_j}$ with respect to a parameter $\theta$ is the
classical identity
\begin{equation}\label{eq:softmax-deriv}
\frac{\partial z_i}{\partial \theta} = z_i \biggl(\frac{\partial x_i}
  {\partial \theta} - \sum_j z_j \frac{\partial x_j}{\partial \theta}\biggr).
\end{equation}
Applying~\eqref{eq:softmax-deriv} with $z_i = \pi^T(a \mid s)$,
$x_i = Q_{R_\varepsilon}(s,a)/T$, and $\theta = \varepsilon$, we obtain
\[
G(s,a) = \pi^T(a \mid s) \biggl(\frac{1}{T}\widetilde{\Phi}(s,a) -
  \sum_{b} \pi^T(b \mid s) \frac{1}{T}\widetilde{\Phi}(s,b)\biggr),
\]
which is~\eqref{eq:sensitivity}.
\end{proof}

\begin{remark}\label{rem:phi-tilde}
The Q-function sensitivity $\widetilde{\Phi}$ incorporates both the direct
effect of the perturbation $\Phi$ at the current state and its discounted
propagation through future states. In many formulas below, we write $\Phi$ in
place of $\widetilde{\Phi}$ for notational clarity, with the understanding
that all statements hold for the effective perturbation direction.
\end{remark}

We now state and prove the main characterization theorem.

\begin{theorem}[Threshold characterization]\label{thm:characterization}
Under the conditions of Theorem~\ref{thm:existence}, the misalignment
sensitivity threshold satisfies
\begin{equation}\label{eq:threshold-formula}
\tau(\delta) = \sqrt{\frac{2\delta}{\lVert G \rVert_F^2}} + O(\delta)
\end{equation}
as $\delta \to 0$, where $\lVert G \rVert_F^2$ is the Fisher information
norm of the policy gradient sensitivity.
\end{theorem}

\begin{proof}
The proof proceeds in three steps: we expand $D(\varepsilon)$ to second order,
identify the leading term with the Fisher information, and invert the relation.

\emph{Step 1: Taylor expansion of the per-state KL divergence.}
Fix a state $s \in \mathcal{S}$ and write $\pi_\varepsilon = \pi_\varepsilon
(\cdot \mid s)$, $\pi_0 = \pi_0(\cdot \mid s)$, $G = G(s,\cdot)$ for
brevity. The first-order expansion gives
\[
\pi_\varepsilon(a \mid s) = \pi_0(a \mid s) + \varepsilon \, G(s,a)
  + O(\varepsilon^2).
\]
The per-state KL divergence is
\[
D_{\mathrm{KL}}(\pi_\varepsilon(\cdot \mid s) \| \pi_0(\cdot \mid s))
  = \sum_a \pi_\varepsilon(a \mid s) \log \frac{\pi_\varepsilon(a \mid s)}
  {\pi_0(a \mid s)}.
\]
Substituting $\pi_\varepsilon = \pi_0 + \varepsilon G + O(\varepsilon^2)$ and
using $\log(1+x) = x - x^2/2 + O(x^3)$ with $x = \varepsilon G(s,a) /
\pi_0(a \mid s)$, we obtain
\begin{align}
\sum_a \pi_\varepsilon(a \mid s) \log \frac{\pi_\varepsilon(a \mid s)}
  {\pi_0(a \mid s)}
&= \sum_a \bigl(\pi_0 + \varepsilon G\bigr) \biggl(\frac{\varepsilon G}
  {\pi_0} - \frac{\varepsilon^2 G^2}{2\pi_0^2} + O(\varepsilon^3)\biggr)
  \notag \\
&= \varepsilon \sum_a G(s,a) + \frac{\varepsilon^2}{2} \sum_a
  \frac{G(s,a)^2}{\pi_0(a \mid s)} + O(\varepsilon^3). \label{eq:kl-expand}
\end{align}

\emph{Step 2: The first-order term vanishes.}
Since $\pi_\varepsilon(\cdot \mid s)$ is a probability distribution for every
$\varepsilon$, differentiating $\sum_a \pi_\varepsilon(a \mid s) = 1$ gives
$\sum_a G(s,a) = 0$. Therefore the first-order term in~\eqref{eq:kl-expand}
vanishes, and
\begin{equation}\label{eq:kl-second-order}
D_{\mathrm{KL}}(\pi_\varepsilon(\cdot \mid s) \| \pi_0(\cdot \mid s))
  = \frac{\varepsilon^2}{2} \sum_a \frac{G(s,a)^2}{\pi_0(a \mid s)}
  + O(\varepsilon^3).
\end{equation}

\emph{Step 3: Averaging and inversion.}
Averaging~\eqref{eq:kl-second-order} over the stationary distribution and
noting that $d^{\pi_\varepsilon}(s) = d^{\pi_0}(s) + O(\varepsilon)$, we
obtain
\[
D(\varepsilon) = \frac{\varepsilon^2}{2} \lVert G \rVert_F^2
  + O(\varepsilon^3).
\]
Setting $D(\varepsilon) = \delta$ and solving for $\varepsilon$, we find
$\varepsilon^2 = 2\delta / \lVert G \rVert_F^2 + O(\delta^{3/2})$, whence
\[
\tau(\delta) = \sqrt{\frac{2\delta}{\lVert G \rVert_F^2}} + O(\delta).
\qedhere
\]
\end{proof}

\begin{example}\label{ex:two-state}
Consider a two-state, two-action MDP with $\mathcal{S} = \{s_1, s_2\}$,
$\mathcal{A} = \{a_1, a_2\}$, Q-values $Q(s_1, a_1) = 1$, $Q(s_1, a_2) = 0$,
$Q(s_2, a_1) = 0$, $Q(s_2, a_2) = 1$, perturbation $\Phi(s,a) = Q(s,a)$,
uniform stationary distribution $d(s) = 1/2$, and temperature $T$. Then
\[
\pi^T(a_1 \mid s_1) = \frac{e^{1/T}}{e^{1/T} + 1}, \qquad
G(s_1, a_1) = \frac{1}{T} \cdot \frac{e^{1/T}}{(e^{1/T}+1)^2}.
\]
The Fisher information norm is
$\lVert G \rVert_F^2 = \frac{1}{T^2} \cdot \frac{e^{1/T}}{(e^{1/T}+1)^2}$.
For $T = 1$, this gives $\lVert G \rVert_F^2 \approx 0.197$ and
$\tau(0.01) \approx \sqrt{0.02/0.197} \approx 0.319$.
\end{example}

\subsection{Temperature scaling}

\begin{corollary}[Temperature scaling law]\label{cor:temperature}
The Fisher information norm scales as $\lVert G \rVert_F^2 \propto 1/T^2$,
and consequently
\begin{equation}\label{eq:temp-scaling}
\tau(\delta) \propto T.
\end{equation}
\end{corollary}

\begin{proof}
From Lemma~\ref{lem:sensitivity},
\[
G(s,a) = \frac{1}{T}\, \pi^T(a \mid s) \bigl(\Phi(s,a) -
  \mathbb{E}_{\pi^T}[\Phi(s,\cdot)]\bigr).
\]
Substituting into the Fisher information norm~\eqref{eq:fisher}:
\begin{align}
\lVert G \rVert_F^2 &= \sum_s d(s) \sum_a \frac{1}{\pi_0(a \mid s)}
  \cdot \frac{1}{T^2} \pi_0(a \mid s)^2 \bigl(\Phi(s,a) -
  \mathbb{E}_{\pi_0}[\Phi]\bigr)^2 \notag \\
&= \frac{1}{T^2} \sum_s d(s) \sum_a \pi_0(a \mid s) \bigl(\Phi(s,a) -
  \mathbb{E}_{\pi_0}[\Phi]\bigr)^2 \notag \\
&= \frac{1}{T^2} \sum_s d(s)\, \operatorname{Var}_{\pi_0(\cdot \mid s)}
  \bigl(\Phi(s,\cdot)\bigr). \label{eq:fisher-var}
\end{align}
The variance term $V := \sum_s d(s)\, \operatorname{Var}_{\pi_0(\cdot \mid s)}
(\Phi(s,\cdot))$ depends on $T$ through the policy $\pi_0$, but the dominant
scaling is captured by the explicit $1/T^2$ prefactor. We verify this in two
regimes.

\emph{High temperature} ($T \gg \Delta Q$, where $\Delta Q = \max_{s,a}
Q(s,a) - \min_{s,a} Q(s,a)$): the policy approaches uniform, so $V \to
\operatorname{Var}_{\mathrm{unif}}(\Phi)$, which is a positive constant
independent of $T$. Thus $\lVert G \rVert_F^2 \sim V_\infty / T^2$.

\emph{Low temperature} ($T \ll \Delta Q$): the policy concentrates on the
greedy action, and $\operatorname{Var}_{\pi_0}(\Phi)$ depends on how
$\Phi$ values compare at the greedy action versus others. For generic $\Phi$,
$V$ remains bounded and bounded away from zero. The leading-order behavior
$\lVert G \rVert_F^2 \sim C/T^2$ with $C = C(Q, \Phi) > 0$ persists.

In both regimes, $\lVert G \rVert_F^2 = C(Q,\Phi)/T^2$ at leading order.
By Theorem~\ref{thm:characterization},
\[
\tau(\delta) = \sqrt{\frac{2\delta}{\lVert G \rVert_F^2}}
  = T \sqrt{\frac{2\delta}{C(Q,\Phi)}} \propto T.
\qedhere
\]
\end{proof}

\begin{remark}\label{rem:tradeoff}
Corollary~\ref{cor:temperature} reveals a fundamental
\emph{robustness--capability tradeoff}: increasing $T$ raises the threshold
$\tau$ (more robust to misalignment) but also increases the entropy of the
policy (less decisive, lower expected reward). This tradeoff is intrinsic to
softmax parameterization and cannot be circumvented within this class of
policies.
\end{remark}

\subsection{Weak dependence on market parameters}

We now show that the market parameters $\sigma$ (volatility) and $h$
(decision horizon) enter the threshold only at subleading order.

\begin{proposition}[Weak volatility dependence]\label{prop:volatility}
For a trading MDP with Gaussian transition kernel~\eqref{eq:gaussian-kernel}
of width $\sigma$, the threshold satisfies
\begin{equation}\label{eq:vol-bound}
\tau(\sigma) = \tau_0 + O(\sigma^{-2})
\end{equation}
as $\sigma \to \infty$, where $\tau_0$ is the threshold in the
uniform-transition limit.
\end{proposition}

\begin{proof}
The dependence of $\lVert G \rVert_F^2$ on $\sigma$ enters through the
Q-values $Q_{R}(s,a)$, which are determined by the Bellman equation
involving $\mathcal{T}_\sigma$, and through the stationary distribution
$d^{\pi_0}$.

For the Gaussian kernel~\eqref{eq:gaussian-kernel}, as $\sigma \to \infty$
the transition probabilities converge to the uniform distribution:
$\mathcal{T}_\sigma(s' \mid s,a) \to 1/|\mathcal{S}|$. The rate of
convergence is determined by the Gaussian tail: for fixed $s, s', a$,
\[
\mathcal{T}_\sigma(s' \mid s, a) = \frac{1}{|\mathcal{S}|} + O(\sigma^{-2})
\]
since the exponential weights $\exp(-(s_j - \mu)^2/(2\sigma^2))$ expand as
$1 - (s_j - \mu)^2/(2\sigma^2) + O(\sigma^{-4})$.

The Bellman operator is a contraction with modulus $\gamma < 1$, and the
Q-function depends Lipschitz-continuously on the transition kernel with
constant $R_{\max}/(1-\gamma)^2$, where $R_{\max} = \max_{s,a} |R(s,a)|$.
Therefore
\[
Q_\sigma(s,a) - Q_\infty(s,a) = O(\sigma^{-2}).
\]
Since the Fisher information norm $\lVert G \rVert_F^2$ depends on
$Q_\sigma$ through the softmax policy, which is a smooth function of the
Q-values, we conclude
\[
\lVert G \rVert_F^2(\sigma) = \lVert G \rVert_F^2(\infty) + O(\sigma^{-2}),
\]
and therefore $\tau(\sigma) = \tau_\infty + O(\sigma^{-2})$.
\end{proof}

\begin{proposition}[Weak horizon dependence]\label{prop:horizon}
For effective horizon $h = 1/(1-\gamma)$, the threshold satisfies
\begin{equation}\label{eq:horizon-bound}
\tau(h) = \tau_0 + O(h^{-1})
\end{equation}
as $h \to \infty$.
\end{proposition}

\begin{proof}
The discount factor is $\gamma = 1 - 1/h$. The Q-function satisfies
\[
Q(s,a) = R(s,a) + \gamma \sum_{s'} \mathcal{T}_\sigma(s' \mid s, a) V(s'),
\]
where $V(s) = \max_a Q(s,a)$. As $\gamma \to 1$ (equivalently $h \to
\infty$), the Q-values grow as $Q \sim R_{\mathrm{avg}} \cdot h$, but the
Q-value \emph{differences} $Q(s,a) - Q(s,b)$---which are the only quantities
entering the softmax policy (Remark~\ref{rem:softmax-invariance})---converge
to a finite limit. Indeed,
\[
Q(s,a) - Q(s,b) = R(s,a) - R(s,b) + \gamma \sum_{s'} \bigl[\mathcal{T}(s'
  \mid s, a) - \mathcal{T}(s' \mid s, b)\bigr] V(s'),
\]
and the value differences $V(s') - V(s'')$ remain bounded by
$R_{\max}/(1-\gamma)$ but cancel in the sum structure, yielding convergence
of Q-differences at rate $O(\gamma) = O(1 - 1/h)$.

Since the softmax policy depends only on Q-differences (which converge as
$h \to \infty$), the Fisher information norm converges:
\[
\lVert G \rVert_F^2(h) = \lVert G \rVert_F^2(\infty) + O(h^{-1}),
\]
and therefore $\tau(h) = \tau_\infty + O(h^{-1})$.
\end{proof}

\begin{example}\label{ex:volatility}
Consider the two-state MDP of Example~\ref{ex:two-state} with a Gaussian
transition kernel of width $\sigma$. For $T = 0.3$ and $\delta = 0.01$, the
theoretical threshold is $\tau \approx 0.14$ for all $\sigma \in [0.05, 50]$,
with a variation of less than $2\%$ across three orders of magnitude in
$\sigma$. This illustrates the content of Proposition~\ref{prop:volatility}:
the volatility parameter is essentially irrelevant to misalignment
vulnerability.
\end{example}

\subsection{Computational verification}\label{sec:computation}

We verify the theoretical predictions on finite trading MDPs with state space
of size $n_{\mathrm{prices}} \times n_{\mathrm{positions}}$, action space
$\mathcal{A} = \{\text{sell}, \text{hold}, \text{buy}\}$, Gaussian transition
kernels, and softmax policies. The threshold is computed numerically via
bisection over random perturbation directions. Theoretical predictions use
$\tau_{\mathrm{theory}} = \sqrt{2\delta / F_{\max}}$ where $F_{\max}$ is the
Fisher information evaluated at the worst-case perturbation direction.

\emph{Temperature dependence.}
Table~\ref{tab:temperature} reports thresholds for $T \in [0.05, 3.0]$.
A power-law fit yields $\tau_{\mathrm{numerical}} \propto T^{1.008}$ with
$R^2 = 0.9992$ and $p < 10^{-13}$, confirming the linear scaling of
Corollary~\ref{cor:temperature}. The theoretical prediction
$\tau_{\mathrm{theory}} \propto T^{0.987}$ has $R^2 = 0.9993$. Numerical
and theoretical values agree within $10\%$ across two orders of magnitude
in $T$.

\begin{table}[ht]
\centering
\caption{Temperature dependence of the misalignment threshold. Power-law
  fit: $\tau \propto T^{1.008}$, $R^2 = 0.9992$.}\label{tab:temperature}
\begin{tabular}{@{}cccc@{}}
\toprule
$T$ & $\tau$ (numerical) & $\tau$ (theoretical) & $\lVert G \rVert_F^2$ \\
\midrule
0.05 & 0.0219 & 0.0215 & 43.24 \\
0.10 & 0.0431 & 0.0466 & 9.21 \\
0.20 & 0.0803 & 0.0879 & 2.59 \\
0.50 & 0.2148 & 0.2319 & 0.37 \\
1.00 & 0.4260 & 0.4224 & 0.11 \\
2.00 & 0.8718 & 0.8797 & 0.03 \\
3.00 & 1.3231 & 1.2459 & 0.01 \\
\bottomrule
\end{tabular}
\end{table}

\emph{Volatility dependence.}
The threshold is approximately constant ($\tau \approx 0.14$) across three
orders of magnitude in $\sigma \in [0.05, 50]$. A power-law fit gives
$\tau \propto \sigma^{0.016}$ with $R^2 = 0.43$ (poor fit), confirming
Proposition~\ref{prop:volatility}.

\emph{Horizon dependence.}
The threshold is approximately constant ($\tau \approx 0.21$) for
$h \in \{3, 10, 30, 100\}$. A power-law fit gives
$\tau \propto h^{-0.009}$ with $R^2 = 0.13$ (poor fit), confirming
Proposition~\ref{prop:horizon}.

\emph{Architecture dependence.}
Table~\ref{tab:architecture} reports the exponents $\alpha, \beta$ (from
fits $\tau \propto \sigma^{-\alpha} h^{-\beta}$) across five MDP
architectures. All exponents are within $0.04$ of zero, regardless of
state-space size, temperature, or spectral gap of the transition matrix.

\begin{table}[ht]
\centering
\caption{Estimated exponents $\alpha$ and $\beta$ across MDP architectures.
  All values are negligibly small.}\label{tab:architecture}
\begin{tabular}{@{}lccccc@{}}
\toprule
Architecture & $|\mathcal{S}|$ & $T$ & $\alpha$ & $\beta$ &
  Spectral gap \\
\midrule
Tiny   &  24 & 0.3 & $\phantom{-}0.023$ & $\phantom{-}0.039$ & 0.088 \\
Small  &  48 & 0.5 & $-0.018$            & $-0.021$            & 0.038 \\
Medium & 100 & 0.5 & $\phantom{-}0.018$  & $\phantom{-}0.013$ & 0.013 \\
Cold   & 100 & 0.1 & $-0.021$            & $\phantom{-}0.004$ & 0.013 \\
Warm   & 100 & 1.0 & $\phantom{-}0.027$  & $-0.021$            & 0.013 \\
\bottomrule
\end{tabular}
\end{table}
